\documentclass[11pt]{article}
\usepackage{progtable}

\begin{document}
\progheader{Cinquième}{5e}

\begin{proglist}

\seq{5C01}{Nombres relatifs — introduction}{NUM}{Introduire l’ensemble ℤ et la notion de signe}{
\item 5C01-NUM01 — Comprendre et situer les nombres relatifs
\item 5C01-NUM02 — Comparer et ordonner des relatifs
\item 5C01-NUM03 — Représenter sur une droite graduée
}

\seq{5C02}{Addition et soustraction de relatifs}{CALC}{Maîtriser additions et soustractions sur ℤ}{
\item 5C02-CALC01 — Additionner des nombres relatifs
\item 5C02-CALC02 — Soustraire des nombres relatifs
\item 5C02-CALC03 — Simplifier une somme algébrique
}

\seq{5C03}{Géométrie — triangles et angles}{GEO}{Approfondir la géométrie des triangles}{
\item 5C03-GEO01 — Utiliser les propriétés des triangles
\item 5C03-GEO02 — Calculer des angles (droites transversales)
\item 5C03-GEO03 — Construire la médiatrice et la bissectrice
}

\seq{5C04}{Fractions — opérations}{NUM}{Calculer avec les fractions}{
\item 5C04-NUM04 — Additionner des fractions (même dénom.)
\item 5C04-NUM05 — Réduire au même dénominateur
\item 5C04-NUM06 — Multiplier et diviser des fractions
}

\seq{5C05}{Aires et périmètres — figures composées}{MEAS}{Calculer sur des figures composées}{
\item 5C05-MEAS01 — Calculer aire et périmètre du disque
\item 5C05-MEAS02 — Décomposer des figures complexes
\item 5C05-MEAS03 — Distinguer périmètre et aire
}

\seq{5C06}{Probabilités et statistiques}{STAT}{Initier à la pensée statistique}{
\item 5C06-STAT01 — Calculer une fréquence
\item 5C06-STAT02 — Lire et construire un tableau statistique
\item 5C06-STAT03 — Introduction à la notion de probabilité
}

\end{proglist}
\end{document}
